%% notes-tex/wet-lab/w037-isobutanol-scrnaseq/sections/2-on-hand.tex
%% [[experiments.W037-isobutanol-scrnaseq.strain-selection]]
%% https://github.com/Mjvolk3/torchcell/tree/main/notes-tex/wet-lab/w037-isobutanol-scrnaseq/sections/2-on-hand.tex
%%
%% SOURCE: experiments/W037-isobutanol-scrnaseq/scripts/strain_tables.py
%% Table 1 is generated by that script from the Strain records it holds. Every
%% record carries the Supplementary Table 1 row or main-text sentence it was read
%% from, plus a verbatim quote. The host genotype quoted in the prose is the
%% CEN.PK2-1C row of that same table.

\section{What is on hand\secstatus{todo}}\label{sec:onhand}

The twelve come from one paper \mirror{zhangBiosensorBranchedchainAmino2022}.
Values below were read from the primary article and its Supplementary
Information, whose sha256 are pinned in the generating script and match the
copies the literature mirror holds under that citation key, so every value here
traces to a hash-pinned artifact rather than to a live URL.

\notesec{The host is CEN.PK2-1C, not BY4741}
Supplementary Table~1 of the source paper gives the parent as
\org{S.\ cerevisiae} CEN.PK2-1C, genotype \texttt{MATa ura3-52 trp1-289
leu2-3,112 his3-1 MAL2-8c SUC2}. No strain in the paper derives from S288C or
BY4741. Two consequences follow. Reads should be mapped against a CEN.PK
assembly, because CEN.PK and S288C differ enough that an S288C-only reference
loses reads and miscalls variants. And the reference needs explicit entries for
the transgenes, since none of them are in any yeast assembly: yEGFP, the three
\gene{LEU4} alleles, \gene{Bs\_alsS}, \gene{Ec\_ilvC}\gsup{P2D1-A1},
\gene{Ll\_ilvD}, \gene{Ll\_adhA}\gsup{RE1}, \gene{EL222}, and the
\texttt{hphMX}, \texttt{kanMX}, \texttt{natMX} and \texttt{ShBle} markers.

The \gene{LEU4} alleles need care beyond being listed. A strain can be
\gene{leu4}$\Delta$ and still carry \gene{LEU4} on a plasmid, and the plasmid
allele may be full length, truncated at residue 410, or deleted for S547. Reads
from a 3\gsup{\textquotesingle}-biased protocol will not separate those alleles
unless the reference is built to separate them, and a \gene{leu4}$\Delta$ strain
carrying a plasmid copy will otherwise appear to express a gene it has lost.

%% GENERATED FILE -- do not hand-edit.
%% SOURCE: experiments/W037-isobutanol-scrnaseq/scripts/strain_tables.py
\begin{table}[H]\centering
\footnotesize
\caption[Strains on hand]{The twelve strains on hand, all of them
\org{Saccharomyces cerevisiae} CEN.PK2-1C derivatives. \emph{Titer} is the
48\,h high-cell-density value for the product named beside it, and the key
below the table says how each value is known.
Carbon source differs across the set and is not comparable between rows: the
three figures these numbers come from use 15\% glucose, 15\% galactose, and
2\% glucose respectively. The last column marks the six selected for the first
run.}\label{tab:inventory}
\begin{tabular}{llllll}
\toprule
id & strain & product & titer (mg/L) & carbon & first run \\
\midrule
JC0043 & YZy311 & isobutanol & $\sim$118\,$^{\dagger}$ & 15\% glucose & \checkmark \\
JC0044 & YZy313 & isobutanol & not reported & 15\% glucose & \checkmark \\
JC0045 & YZy314 & isobutanol & 1.6$\times$ JC0044 & 15\% glucose & \checkmark \\
JC0046 & YZy148 & isopentanol & not reported & 15\% glucose &  \\
JC0047 & YZy148 + LEU4 WT & isopentanol & $\sim$201\,$^{\dagger}$ & 15\% glucose &  \\
JC0048 & YZy148 + LEU4dS547 & isopentanol & \textbf{842 $\pm$ 23} & 15\% glucose &  \\
JC0049 & YZy452 & isobutanol & \textbf{310 $\pm$ 15} & 15\% galactose &  \\
JC0050 & YZy454 & isobutanol & $\sim$73\,$^{\dagger}$ & 15\% glucose & \checkmark \\
JC0051 & YZy469 & isobutanol & \textbf{227 $\pm$ 13} & 15\% glucose & \checkmark \\
JC0052 & YZy91 & isobutanol & not reported & 15\% glucose & \checkmark \\
JC0053 & YZy502 & isobutanol & not reported & 2\% glucose, dark &  \\
JC0054 & SHy134 & isopentanol & not reported & 15\% glucose &  \\
\bottomrule
\multicolumn{6}{@{}l@{}}{\scriptsize \textbf{bold} stated in the paper as a number;\quad $\dagger$ derived by dividing a stated fold change, so approximate and without uncertainty} \\
\end{tabular}
\end{table}


\notesec{Four strains make isopentanol, which splits the set in two}
\strain{JC0046} through \strain{JC0048} and \strain{JC0054} produce isopentanol; the
remaining eight produce isobutanol. There is no single production ranking across
the twelve, and the highest titer on hand, 842\,$\pm$\,23\,mg/L from
\strain{JC0048}, is an isopentanol number.

\notesec{The paper's best producers are not on hand}
The strains that set the paper's production ceiling are absent from the twelve:
YZy505 at 681\,$\pm$\,29\,mg/L isobutanol, YZy148 carrying \gene{LEU4} mutant 6
at 963\,$\pm$\,32\,mg/L isopentanol, YZy470 at 443\,$\pm$\,6\,mg/L and YZy312 at
378\,$\pm$\,10\,mg/L. \strain{JC0053} is YZy502, the plasmid-free parent of
YZy505, and not YZy505 itself. Reaching the paper's ceiling means obtaining or
rebuilding strains beyond this set, which is a separate decision from the one
this run makes.
